% Generated draft: requires mathematical and source review.
\documentclass[11pt]{article}
\usepackage[margin=25mm]{geometry}
\usepackage{fontspec}
\setmainfont{Latin Modern Roman}
\usepackage{amsmath,amssymb,mathtools,unicode-math}
\setmathfont{Latin Modern Math}
\usepackage{enumitem,needspace,xurl,hyperref}
\setlength{\parindent}{0pt}
\setlength{\parskip}{5pt}
\setlength{\emergencystretch}{3em}
\newcommand{\N}{\mathbb{N}}
\newcommand{\Z}{\mathbb{Z}}
\newcommand{\Q}{\mathbb{Q}}
\newcommand{\R}{\mathbb{R}}
\newcommand{\C}{\mathbb{C}}
\newcommand{\F}{\mathbb{F}}
\newcommand{\A}{\mathbb{A}}
\newcommand{\PP}{\mathbb{P}}
\newcommand{\E}{\mathbb{E}}
\newcommand{\eps}{\varepsilon}
\newcommand{\dd}{\,\mathrm{d}}
\begin{document}
\section*{UnsolvedMath ID: 30002854}
\textbf{Source problem code:} catalogue entry
\subsection{Definitions and mathematical statement}
For the definitions below, let \(G\) be a finite group, let \(p\) be a prime dividing \(|G|\), and let \(k\) be a field of characteristic \(p\). This supplies the usual modular setting for interpreting the question; the record does not specify the coefficient field, whether it must be algebraically closed, or explicitly state the divisibility assumption.

A finite-dimensional \(kG\)-module \(M\) is \emph{endotrivial} if
\[
\operatorname{End}_k(M)\cong k\oplus P
\]
as \(kG\)-modules, where \(k\) is the trivial module, \(P\) is projective, and \(G\) acts on endomorphisms by \((g\cdot f)(m)=g f(g^{-1}m)\). Two modules are stably isomorphic if they become isomorphic after adding finite-dimensional projective modules. The endotrivial-module group \(T_k(G)\), written \(T(G)\) in the source, consists of stable isomorphism classes of endotrivial modules, with multiplication induced by tensor product over \(k\) with diagonal \(G\)-action, identity \([k]\), and inverse induced by the dual module. Its torsion subgroup is
\[
TT_k(G)=T_T(G)=\bigl\{[M]\in T_k(G):\text{ there exists }n\geq 1\text{ with }[M^{\otimes_k n}]=[k]\bigr\}.
\]
Here \(T_T(G)\) denotes the source's alternative notation for this subgroup, with dependence on \(k\) suppressed. Define
\[
O_{p'}(G)=\text{the largest normal subgroup of }G\text{ whose order is relatively prime to }p.
\]

\textbf{Source question.} Can the determination of \(TT(G)\) for finite groups be reduced essentially to its determination for groups \(H\) satisfying \(O_{p'}(H)=1\)?

The word ``essentially'' is not defined in the supplied record. No reduction procedure, permitted auxiliary data, reconstruction statement, or comparison map is specified. Consequently, the question above preserves the supplied target but does not specify a fully quantified reduction theorem. In particular, the role of \(G/O_{p'}(G)\), the coefficient-field hypotheses, and the intended scope of uniformity in \(G\) and \(p\) require clarification.
\subsection{Short English statement}
The problem asks whether understanding the finite-order classes of endotrivial modules can be reduced to studying finite groups with no nontrivial normal subgroup of order prime to \(p\). The intended meaning of the reduction and the coefficient-field assumptions are unspecified. The supplied record labels the problem open; this draft does not independently verify that status.
\subsection{Sources}
\begin{itemize}
\item Supplied local UnsolvedMath registry record, permanent ID \(30002854\), source problem code \texttt{catalogue entry}, in the collection \emph{Oberwolfach Reports Open Problems}. No individual proposer is supplied. The proposed year is recorded as \(2015\).
\item Supplied report citation: \emph{Representations of Finite Groups (2016)}, Oberwolfach Reports, pp.\ 971--1027, DOI \texttt{10.4171/owr/2015/18}. Original report link: \url{https://doi.org/10.4171/owr/2015/18}. The supplied title year and the year component of the DOI differ; both are retained without resolving the bibliographic discrepancy.
\item Supplied literature source: Carlson--Mazza--Th\'evenaz, \emph{Endo-trivial modules with dihedral Sylow 2-subgroup}, \url{https://arxiv.org/abs/1507.01406}.
\item Literature assessment dated exactly \texttt{2026-08-21}: the supplied curation labels the problem open and states that available results handle important cases and quotients but do not establish the exact general reduction requested here. This is supplied dated curation, described in the record as literature triage derived from the OWR curation; it is not independently verified. The time-sensitive status and cited claims require independent verification.
\item The registry labels the statement \texttt{corrected\_verified} and says its mathematical content was preserved while notation, typography, or omitted definitions were repaired for publication. This is a supplied editorial designation, not an independent mathematical verification of this draft.
\end{itemize}
\end{document}
